\documentclass[11pt]{article}
\usepackage[margin=1in]{geometry}
\usepackage[T1]{fontenc}
\usepackage[utf8]{inputenc}
\usepackage{lmodern,microtype,amsmath,amssymb,amsthm,mathtools}
\usepackage[colorlinks=true,linkcolor=blue,citecolor=blue,urlcolor=blue]{hyperref}
\hypersetup{pdftitle={Exponential online mistakes for tanh networks with stronger readout regularization},pdfauthor={ChatGPT-6 Astra; Kaarel Hanni; Dmitry Vaintrob}}
\newtheorem{theorem}{Theorem}
\newtheorem{lemma}[theorem]{Lemma}
\newtheorem{proposition}[theorem]{Proposition}
\newtheorem{corollary}[theorem]{Corollary}
\newtheorem{remark}[theorem]{Remark}
\newcommand{\R}{\mathbb R}
\newcommand{\E}{\mathbb E}
\newcommand{\Pp}{\mathbb P}
\newcommand{\sgn}{\operatorname{sign}}
\title{Exponential online mistakes for tanh networks\protect\\with stronger readout regularization}
\author{ChatGPT-6 Astra \and Kaarel H\"anni \and Dmitry Vaintrob\thanks{Authors are listed in alphabetical order.}}
\date{6 November 2026}
\begin{document}
\maketitle

\begin{abstract}
A Bayesian neural network can learn a small target from polynomially
many IID examples and nevertheless make exponentially many mistakes on
one fixed ordering of the input cube. We prove this for ordinary tanh,
one hidden layer, independent Gaussian weights with equal variances
across all first-layer coordinates, and every sufficiently large hidden
width. The prior has first-layer variance $\tau^2N^{1-\gamma}/d$ and
readout variance $\tau^2N^{-\gamma-1}/A_N$, where $\gamma\ge1$, $\tau>0$
are fixed and $A_N\to\infty$ can diverge arbitrarily slowly.
The target has a Boolean circuit of size $\Theta(k)$ and a shallow tanh
teacher of width $O(k^4)$ with every weight bounded by one. After an
$O(k)$ calibration prefix, a fixed input permutation forces at least
$2^{k-1}$ designated mistakes, with posterior wrong-label probability
tending to one as $N\to\infty$. The same posterior has polynomial IID
sample complexity. The proof gives an exact scalar-optimality theorem:
sufficiently strong readout regularization makes the minimum energy of
any width equal to an explicit function of the affine margin norm.
A Gaussian-volume comparison then handles the growing parameter dimension.
A subpolynomial readout correction, such as a logarithmic one, preserves
the corresponding power-law exponents. The result does not establish the
case $A_N\equiv1$ or a lower bound at every realizable width.
\end{abstract}

\section{Model and main result}

An IID guarantee controls errors when training and testing use the same
input distribution. An online mistake bound would be stronger: it would
control the total number of errors even when the inputs arrive in an
unfavorable order. We exhibit a separation between these guarantees for
an exact Bayesian tanh learner. The network and target use the same
shallow architecture; the lower bound persists at arbitrarily large
learner width.

Let $X=\{-1,1\}^d$ and $\bar x=(1,x)\in\R^{D}$, where $D=d+1$.
At hidden width $N$, the learner is
\begin{equation}\label{eq:model}
 F_{u,v}(x)=\sum_{j=1}^N v_j\tanh(u_j\cdot\bar x).
\end{equation}
All weights are independent centered Gaussians, with
\begin{equation}\label{eq:prior}
 \operatorname{Var}(u_{ji})=\frac{\tau^2N^{1-\gamma}}d,
 \qquad
 \operatorname{Var}(v_j)=\frac{\tau^2N^{-\gamma-1}}{A_N}.
\end{equation}
The constant input coordinate has the same first-layer variance as every
other coordinate. Fix $\gamma\ge1$, $\tau>0$, and any positive sequence
$A_N\to\infty$. No monotonicity or minimum rate of divergence is assumed.
For example one may use
$A_N=\log\log(N+e^e)$, or a slower divergent sequence.

For a Boolean function $g$ define its strict global-margin cell
\begin{equation}\label{eq:cell}
 U_g=\{(u,v):g(x)F_{u,v}(x)>1/2\text{ for every }x\in X\}.
\end{equation}
The prior is conditioned on the union of these cells. After every
prediction the true label is revealed, and the posterior retains exactly
the cells that fit all observed labels. The online input order below has no repetitions.
The posterior error probability is the mass of functions with the wrong
label at the next input. This is also the error probability of a fresh
posterior sample; posterior-label majority is wrong when it exceeds
$1/2$. The confidence convention requiring more than $2/3$ posterior
probability on the correct label is covered as well.

The factor $A_N$ is the substantive prior change in this theorem.
The reference normalization has $A_N=1$, with readout variance
$\tau^2N^{-\gamma-1}$. We prove the lower bound after dividing that
variance by any diverging factor. If this factor is subpolynomial,
the layerwise power exponents are unchanged. The result does not remove
the factor altogether.

\begin{theorem}\label{thm:main}
For every $k\ge3$, put $d=2k+1$ and write inputs as $(x,y,z)$ with
$x,y\in\{-1,1\}^k$ and $z\in\{-1,1\}$. Define
\[
 X_x=\sum_{i=0}^{k-1}2^i\frac{x_i+1}{2},\qquad
 Y_y=\sum_{i=0}^{k-1}2^i\frac{y_i+1}{2},
 \qquad
 f_k(x,y,z)=+1\ \Longleftrightarrow\ z=-1\text{ and }X_x\ge Y_y.
\]
There is one fixed permutation of the input cube with the following
properties.
\begin{enumerate}
\item The target has a bounded-fan-in Boolean circuit of size $\Theta(k)$.
It also has a one-hidden-layer tanh teacher with all weights bounded by
one, global output margin greater than $3/4$, and hidden width at most
\[
 n_k=\left\lceil288(k+2)^2(4k+3)(2k+3)\log2\right\rceil=O(k^4).
\]
\item The permutation begins with at most $4k+3$ calibration examples
and then has $T_k=2^k-k-1\ge2^{k-1}$ designated rounds. For every fixed
choice of $\gamma\ge1$, $\tau>0$, and $A_N\to\infty$ in
\eqref{eq:prior}, the posterior wrong-label probability tends to one
at each designated round as $N\to\infty$.
\item Thus at every sufficiently large width, posterior-label majority
makes all $T_k$ mistakes and the strict $2/3$ correctness criterion fails
at all those rounds. For every $\eta>0$, posterior sampling makes all
$T_k$ mistakes with probability greater than $1-\eta$ at every
sufficiently large width.
\item For the same prior and target,
$O((d^2+\log(1/\beta))/\varepsilon)$ IID examples suffice, with
probability at least $1-\beta$, for posterior correct-label probability
greater than $2/3$ on input-distribution mass at least $1-\varepsilon$,
at every sufficiently large width. The width threshold can be independent
of the training sample and its input distribution.
\end{enumerate}
The target and permutation depend only on $k$. The width threshold may
depend on $k,\gamma,\tau,A_N$, and the required confidence. It is not
claimed to be polynomial, or to equal the teacher width.
\end{theorem}

The errors are exponential in the target's Boolean circuit size, and
superpolynomial in the displayed bounded-weight shallow teacher bound.
No choice among tied minimum-energy predictors is used. The comparison
below bounds every network that would make the correct next prediction.

\section{Exact scalar optimality at every width}

Write $b_j=Nv_j$ and temporarily fix a readout penalty $A>0$.
The Gaussian rate energy is
\begin{equation}\label{eq:energy}
 E_A=\frac dN\sum_{j=1}^N\|u_j\|^2+\frac AN\sum_{j=1}^Nb_j^2,
 \qquad F(x)=\frac1N\sum_jb_j\tanh(u_j\cdot\bar x).
\end{equation}
The next theorem concerns arbitrary finite linearly separable training
sets. Its lower bound does not assume that the network's output sign
function is an affine threshold.

Let $z_i=y_i\bar x_i\in\R^D$ be the signed data vectors and let $a_*$
minimize $\|a\|^2$ subject to $z_i\cdot a\ge1$. Put $Q=\|a_*\|^2>0$.
The minimizer is unique. Its KKT conditions give a set $I$ of active
constraints and positive coefficients $\mu_i$ such that
\begin{equation}\label{eq:kkt}
 a_*=\sum_{i\in I}\mu_i z_i,\qquad
 z_i\cdot a_*=1\ (i\in I),\qquad \sum_{i\in I}\mu_i=Q.
\end{equation}
Set $\rho_i=\mu_i/Q$ and
$p=\min\{1/2,\min_{i\in I}\rho_i\}>0$.
Define
\begin{equation}\label{eq:profile}
 \Phi_A(Q)=\min_{r>0}\left\{dQr^2+\frac{A}{4\tanh^2r}\right\}.
\end{equation}
Its unique minimizer $r=r_A(Q)$ satisfies
\begin{equation}\label{eq:stationary}
 A=\frac{4dQr\tanh^3r}{\operatorname{sech}^2r}.
\end{equation}
The right side is strictly increasing from zero to infinity.
In particular $r_A(Q)$ increases with $A$ and decreases with $Q$.
For $Q_1>Q_0$, evaluation at $r_A(Q_1)$ gives
\begin{equation}\label{eq:profile-gap}
 \Phi_A(Q_1)-\Phi_A(Q_0)
 \ge d(Q_1-Q_0)r_A(Q_1)^2>0.
\end{equation}

\begin{theorem}[An exact all-width minimum]\label{thm:scalar}
If
\begin{equation}\label{eq:A-threshold}
 A\ge \frac{256dQ\log(16/p)}{p^2},
\end{equation}
then, for every $N\ge1$, the minimum of \eqref{eq:energy} over all
width-$N$ tanh networks satisfying the training output margins
$y_iF(x_i)\ge1/2$ equals $\Phi_A(Q)$. It is attained by cloning the
single neuron
\[
 u_j=r_A(Q)a_*,\qquad b_j=\frac1{2\tanh r_A(Q)}.
\]
The infimum with strict training margins is the same.
\end{theorem}
\begin{proof}
The displayed clone satisfies every training margin and has energy
$\Phi_A(Q)$. We prove a matching lower bound for every network.
Write $r=r_A(Q)$, and set
\[
 \lambda=\frac A{\tanh^2r},\qquad K=\frac{\lambda^2}{4A},\qquad
 S(u)=\sum_{i\in I}\rho_i\tanh(z_i\cdot u).
\]
For any scalar $b$, completing the square gives
\begin{equation}\label{eq:atom-dual}
 d\|u\|^2+Ab^2-\lambda bS(u)
 \ge d\|u\|^2-KS(u)^2.
\end{equation}
Oddness permits replacing $u$ by $-u$ when necessary, so we may suppose
$S(u)\ge0$ when bounding the right side.

If every active score $z_i\cdot u$ is nonnegative, concavity of tanh on
$[0,\infty)$, \eqref{eq:kkt}, and Cauchy--Schwarz imply
\begin{equation}\label{eq:good-atom}
 S(u)\le\tanh\!\left(\sum_i\rho_i z_i\cdot u\right)
 =\tanh(a_*\cdot u/Q)
 \le\tanh(\|u\|/\sqrt Q).
\end{equation}
Thus \eqref{eq:atom-dual} is at least
\begin{equation}\label{eq:dual-floor}
 \min_{s\ge0}\{dQs^2-K\tanh^2s\}
 =dQr^2-K\tanh^2r=:\zeta.
\end{equation}
To check the equality, $h(v)=\tanh^2\sqrt v$ is strictly concave for
$v>0$. Put $s=\sqrt v$. Then $h'(v)=\tanh(s)\operatorname{sech}^2(s)/s$.
The logarithmic derivative of this expression with respect to $s$ is
\[
 \frac1{\sinh s\cosh s}-2\tanh s-\frac1s<0.
\]
Here $\sinh s\cosh s>s$ for $s>0$.
Consequently the objective in \eqref{eq:dual-floor} is strictly convex
as a function of $s^2$, and \eqref{eq:stationary} makes $s=r$ stationary.

If instead one active score is negative, that term contributes at most
zero while all others contribute at most their weights. Hence
$S(u)\le1-p$, and \eqref{eq:atom-dual} is at least $-K(1-p)^2$.
This is at least $\zeta$ provided
\begin{equation}\label{eq:saturation-conditions}
 \tanh^2r\ge1-p,\qquad Ap/8\ge dQr^2.
\end{equation}
Indeed $\tanh^2r-(1-p)^2\ge p(1-p)\ge p/2$ and $K\ge A/4$.

We verify \eqref{eq:saturation-conditions} from \eqref{eq:A-threshold}.
Put $r_0=\log(16/p)$. Since
$\sinh^3r/\cosh r\le e^{2r}/4$, the right side of
\eqref{eq:stationary} at $r_0$ is at most
$256dQ\log(16/p)/p^2$. Thus $r\ge r_0$.
Then $\operatorname{sech}^2r\le4e^{-2r}\le p^2/64\le p$.
Also $\tanh r\ge1/2$, so \eqref{eq:stationary} gives
\[
 \frac{Ap}{8dQr^2}
 =\frac{p\tanh^3r}{2r\operatorname{sech}^2r}
 \ge\frac{pe^{2r}}{64r}
 \ge\frac4{p\log(16/p)}>1.
\]
The middle inequality uses that $e^{2r}/r$ increases for $r\ge r_0$;
the last uses $p\log(16/p)\le\tfrac12\log32<4$.

We have proved the same atomwise lower bound $\zeta$ in both cases.
Average \eqref{eq:atom-dual} over the neurons. The training constraints
imply $N^{-1}\sum_j b_jS(u_j)\ge1/2$. Therefore
\[
 E_A\ge\lambda/2+\zeta
 =dQr^2+\frac A{4\tanh^2r}=\Phi_A(Q).
\]
Finally, an arbitrarily small readout amplification makes the cloned
training margins strict and approaches the same energy.
\end{proof}

\begin{corollary}[A common Boolean-data threshold]\label{cor:uniform-A}
Suppose $z_i\in\{-1,1\}^D$ and $Q\le Q_{\max}$, where $Q_{\max}\ge1$.
For every such training set, Theorem~\ref{thm:scalar} holds whenever
\begin{equation}\label{eq:A-uniform}
 A\ge A_0:=1024dQ_{\max}^3D^{2D}
                  \log(32Q_{\max}D^D).
\end{equation}
This bound is independent of the number of training examples and of $N$.
\end{corollary}
\begin{proof}
Choose a KKT certificate with minimal positive support. Its active
vectors are linearly independent: otherwise a linear dependence lets
one vary the nonnegative coefficients until one vanishes, preserving
their conic combination. Thus the support size is $r\le D$.
Its Gram matrix $G$ is an integer positive-definite matrix with diagonal
$D$, and $G\mu=\mathbf1$. Hadamard's inequality gives
$1\le\det G\le D^r\le D^D$. Cramer's rule and positivity imply
$\mu_i\ge1/\det G\ge D^{-D}$, since each numerator is a positive
integer. Hence $\rho_i\ge1/(Q_{\max}D^D)$.
The common value $p_0=1/(2Q_{\max}D^D)$ is a valid smaller bound for $p$.
Substituting it and $Q_{\max}$ into \eqref{eq:A-threshold} gives
\eqref{eq:A-uniform}.
\end{proof}

For $D=O(k)$ and $Q_{\max}=\exp(O(k))$, this sufficient penalty is
$\exp(O(k\log k))$. The theorem nevertheless allows any divergent
sequence $A_N$: it eventually exceeds the fixed threshold associated
with each fixed target. A slow choice can make the required width large.

\section{A fixed input order with exponentially many affine mistakes}

We give the complete affine construction used by the network theorem.
Keep the target $f_k$ from Theorem~\ref{thm:main}, and abbreviate the two
encoded integers by $X$ and $Y$. The calibration examples all have
selector $z=-1$.
\begin{enumerate}
\item Take diagonal inputs $x=y$ equal to the all-positive vector,
each of its $k$ one-coordinate flips, and the negatives of these vectors.
Every label is positive.
\item For each $i\in\{0,\ldots,k-1\}$, take two negative examples with
$x_i=-1,y_i=+1$ and $x_j=+1,y_j=-1$ for $j<i$. For $j>i$, set the equal
bits $x_j=y_j$ all positive in one example and all negative in the other.
Their highest differing bit is $i$, so $X<Y$.
\end{enumerate}
Remove duplicate examples. The resulting calibration has at most $4k+2$
inputs. Let $q_t$ denote the input $z=+1$, $Y=0$, $X=t$; its target label
is negative for every $0\le t\le2^k-1$.

\begin{lemma}\label{lem:selector}
For $k\le t\le2^k-1$, the unique minimum-norm affine unit-margin
coefficient vector fitting the calibration and the negative example
$q_t$ is the vector of the affine score
\begin{equation}\label{eq:selector-score}
 a_t\cdot\bar x=-t-(t+1)z+\sum_{i=0}^{k-1}2^i(x_i-y_i).
\end{equation}
It fits all earlier $q_j$ with $j\le t$, and its absolute score is at
least one on the entire cube. Its squared norm is
\begin{equation}\label{eq:Qt}
 Q_t=t^2+(t+1)^2+2\sum_{i=0}^{k-1}4^i.
\end{equation}
For $t<2^k-1$, it predicts $q_{t+1}$ positively, and adding the correct
negative label raises the convex affine minimum by
\begin{equation}\label{eq:Qgap}
 Q_{t+1}-Q_t=4(t+1)>0.
\end{equation}
\end{lemma}
\begin{proof}
Write an arbitrary affine score as
\[
 b+cz+\sum_i a_i x_i+\sum_i h_i y_i,
\]
and set
\[
 u=b-c,\quad v=b+c,\quad
 w_i=(a_i-h_i)/2,\quad s_i=a_i+h_i,\quad U=u-1.
\]
The squared coefficient norm is
\begin{equation}\label{eq:affine-norm}
 \frac{u^2+v^2}{2}+2\sum_iw_i^2+\frac12\sum_i s_i^2.
\end{equation}
The paired positive diagonal constraints give
\[
 U\ge0,\qquad \left|\sum_j s_j\right|\le U,\qquad
 \left|\sum_j s_j-2s_i\right|\le U.
\]
Subtracting the expressions inside the last two absolute values shows
$|s_i|\le U$, hence $\sum_i|s_i|\le kU$.
Average the two negative carry constraints at coordinate $i$. The upper
equal-bit terms cancel and leave
\[
 u-2w_i+2\sum_{j<i}w_j\le-1.
\]
Induction yields
\begin{equation}\label{eq:binary-lower}
 w_i\ge\sum_{j<i}w_j+\frac{u+1}{2}
       \ge2^i(1+U/2).
\end{equation}
At $q_t$, the non-selector part of the input contribution is
$2\sum_{\operatorname{bit}_i(t)=1}w_i-
\sum_{\operatorname{bit}_i(t)=0}s_i$. Its negative unit-margin constraint
therefore implies
\[
 v\le-1-2\sum_{\operatorname{bit}_i(t)=1}w_i
            +\sum_{\operatorname{bit}_i(t)=0}s_i
   \le-(2t+1)+(k-t)U\le-(2t+1),
\]
since $t\ge k$. Combining this with $u\ge1$ and
\eqref{eq:binary-lower} in \eqref{eq:affine-norm} gives
\[
 \|a\|^2\ge\frac{1+(2t+1)^2}{2}+2\sum_i4^i=Q_t.
\]
Equality requires $u=1$, $s_i=0$, $w_i=2^i$ and $v=-(2t+1)$,
which is exactly \eqref{eq:selector-score}.

Conversely the displayed score equals $1+2(X-Y)$ on $z=-1$, and
$-(2t+1)+2(X-Y)$ on $z=+1$. It fits every calibration label, and its score
at $q_j$ is $2j-2t-1\le-1$ for $j\le t$. Both branch scores are odd
integers on the full cube, so the absolute margin is at least one.
At $q_{t+1}$ the score is $+1$. Applying the same minimum calculation
to the augmented history gives $Q_{t+1}$, and subtraction gives
\eqref{eq:Qgap}.
\end{proof}

The input permutation starts with the calibration, then $q_k$, then
$q_{k+1},\ldots,q_{2^k-1}$, and finally every remaining input in
lexicographic order. These blocks are disjoint. Its calibration prefix
has length at most $4k+3$, and there are
$T_k=2^k-k-1\ge2^{k-1}$ designated queries. Just before each one, the
wrong affine witness has the exact convex margin minimum for all data
seen so far. The upper bound $Q_t\le3\cdot4^k$ suffices in
Corollary~\ref{cor:uniform-A}.

The target has a linear-size Boolean circuit. With $0/1$ bit notation,
comparison of the lowest $i+1$ bits obeys
\[
 G_{-1}=1,\qquad
 G_i=(X_i\wedge\neg Y_i)\ \vee\
       ((X_i\leftrightarrow Y_i)\wedge G_{i-1}).
\]
Each step uses a constant number of bounded-fan-in gates; conjoin the
last result with $z=-1$. Every input bit is essential: for $x_i$ use
$Y=2^i$ and all other $X$ bits zero; for $y_i$ use $X=0$ and all other
$Y$ bits zero; for $z$ take $X=Y$. A fan-in-two circuit with $d$ essential
inputs has at least $d-1$ gates, by counting edges in the connected
ancestry of its output. This proves the matching $\Theta(k)$ bound.

\section{Bounded-weight shallow teachers from Fourier coefficients}
\label{sec:teacher}

The target used above has a small teacher at the same shallow depth as
 the learner.  We prove this using a general construction that may also be
 useful for other targets.  All networks in this section have one tanh
 hidden layer, one constant input coordinate, and a linear readout.

For $S\subseteq\{1,\ldots,d\}$, write
$\chi_S(x)=\prod_{i\in S}x_i$.  A Boolean function has the Fourier expansion
\[
 f(x)=\sum_{S\subseteq\{1,\ldots,d\}}\widehat f(S)\chi_S(x),
 \qquad
 \mathcal A(f)=\sum_S|\widehat f(S)|.
\]
The Fourier coefficients here use the uniform measure on the Boolean
cube; the learning theorem does not require a uniform input distribution.

\begin{theorem}[A bounded-weight shallow teacher]\label{thm:fourier-teacher}
Every $f:\{-1,+1\}^d\to\{-1,+1\}$ has a one-hidden-layer tanh
implementation that uniformly approximates $f$ within $1/4$, has every
weight of absolute value at most $1$, and has hidden width at most
\begin{equation}\label{eq:teacher-width}
 \left\lceil
 288\mathcal A(f)^2(2d+1)(d+2)\log 2
 \right\rceil.
\end{equation}
In particular, its signed output is greater than $3/4$ at every input.
\end{theorem}

\begin{lemma}[Parity correlation of bounded tanh features]
\label{lem:parity-feature}
For $r\ge0$, let $b_r=0$ when $r$ is odd and $b_r=1$ when $r$ is even.
For independent uniform signs $s_1,\ldots,s_r$, define
\[
 c_r=\mathbb E_s\left[
    \left(\prod_{i=1}^r s_i\right)
    \tanh\left(\sum_{i=1}^r s_i+b_r\right)
 \right].
\]
Then
\begin{equation}\label{eq:parity-correlation}
 |c_r|\ge\frac{1}{3\sqrt{2r+1}}.
\end{equation}
The empty product and sum give $c_0=\tanh 1$.
\end{lemma}
\begin{proof}
Put $m=\lfloor r/2\rfloor$ and
$\beta_m=\binom{2m}{m}/4^m$.  The highest-degree Fourier coefficient of
$\operatorname{sign}(\sum_i s_i+b_r)$ is $(-1)^m\beta_m$.
For $r\ge1$, this follows by splitting the binomial sum at
$\lfloor r/2\rfloor$ and using
\[
 \sum_{j=0}^{q}(-1)^j\binom rj
   =(-1)^q\binom{r-1}{q};
\]
the case $r=0$ is immediate.

Every possible value of $T=\sum_i s_i+b_r$ is an odd integer, and each
individual value has probability at most $\beta_m$.  For even $r$ this
is the central-binomial bound.  For odd $r=2m+1$, the largest probability
is $\frac{2m+1}{2m+2}\beta_m$.  Since
$|\operatorname{sign}(t)-\tanh t|\le2e^{-2|t|}$ for $t\ne0$,
\begin{align*}
 \left|c_r-(-1)^m\beta_m\right|
 &\le\mathbb E\left|
     \operatorname{sign}(T)-\tanh T\right|\\
 &\le4\beta_m\sum_{j\ge0}e^{-2(2j+1)}
   =\frac{4e^{-2}}{1-e^{-4}}\beta_m
   <\frac{7}{12}\beta_m.
\end{align*}
The last inequality follows from $e^2>7$ and
$4e^{-2}/(1-e^{-4})=4e^2/(e^4-1)<7/12$.
Thus $|c_r|>\frac{5}{12}\beta_m>\frac13\beta_m$.

Finally, $\beta_m\ge1/\sqrt{4m+1}$.  This follows by induction from
$\beta_0=1$ and
$\beta_{m+1}/\beta_m=(2m+1)/(2m+2)$: after squaring, the induction step is
\[
 (2m+1)^2(4m+5)\ge4(m+1)^2(4m+1),
\]
whose left-hand side exceeds the right-hand side by $1$.
Since $4m+1\le2r+1$, the claimed bound follows.
\end{proof}

\begin{proof}[Proof of Theorem~\ref{thm:fourier-teacher}]
For a subset $S$ of cardinality $r$ and a sign vector $s$ indexed by $S$,
let
\[
 h_{S,s}(x)=\tanh\left(\sum_{i\in S}s_i x_i+b_r\right).
\]
Every affine coefficient belongs to $\{0,-1,+1\}$, including the
coefficient of the constant input.  Changing variables from $s_i$ to
$s_i x_i$ gives the exact identity
\begin{equation}\label{eq:parity-mixture}
 \mathbb E_s\left[
   \left(\prod_{i\in S}s_i\right)h_{S,s}(x)
 \right]=c_r\chi_S(x).
\end{equation}
Consequently, the Fourier expansion expresses $f$ as
\[
 f(x)=B\mathbb E[H(x)],
 \qquad
 B=\sum_S\frac{|\widehat f(S)|}{|c_{|S|}|}
   \le3\mathcal A(f)\sqrt{2d+1},
\]
where $H$ is a signed feature of the form $h_{S,s}$ and $|H(x)|\le1$.
Explicitly, choose $S$ with probability proportional to
$|\widehat f(S)|/|c_{|S|}|$, choose $s$ uniformly, and multiply
$h_{S,s}$ by the signs prescribed by
$\widehat f(S)/c_{|S|}$ and $\prod_{i\in S}s_i$.
The empty subset uses the constant feature $\tanh1$ and needs no separate
construction.

Take
\[
 n=\left\lceil32B^2\log(4\cdot2^d)\right\rceil
\]
independent features $H_1,\ldots,H_n$.  Hoeffding's inequality and a union
bound over the cube give
\[
 \mathbb P\left\{
  \max_x\left|\frac Bn\sum_{j=1}^nH_j(x)-f(x)\right|\ge\frac14
 \right\}
 \le2\cdot2^d\exp\left(-\frac{n}{32B^2}\right)
 \le\frac12.
\]
A choice of features therefore achieves uniform error strictly below
$1/4$.  Put each feature's sign in its readout coefficient, whose
magnitude is $B/n$.  Because $\mathcal A(f)\ge1$ and $|c_r|\le1$, we have
$B\ge1$ and $n\ge B$; hence every readout weight also has magnitude at
most $1$.  Finally,
$B^2\le9\mathcal A(f)^2(2d+1)$ and
$\log(4\cdot2^d)=(d+2)\log2$ give~\eqref{eq:teacher-width}.
\end{proof}

\begin{corollary}[A shallow teacher for the selector target]
\label{cor:selector-teacher}
For the target $f_k$ on $d=2k+1$ inputs, equal to $+1$ precisely when
$z=-1$ and $X\ge Y$, there is a one-hidden-layer tanh teacher with all
weights bounded by $1$, global signed output greater than $3/4$, and
width at most
\begin{equation}\label{eq:selector-teacher-width}
 W_k=\left\lceil
 288(k+2)^2(4k+3)(2k+3)\log2
 \right\rceil=O(k^4).
\end{equation}
In the convention that $C_2$ bounds both hidden width and individual
weight magnitudes, $C_2(f_k)\le W_k$.
\end{corollary}
\begin{proof}
Let $x_i,y_i$ be the two integers' Boolean bits, with $i=0$ the least
significant position, and put $E_i=(1+x_i y_i)/2$.  The comparison sign
$G$, equal to $+1$ when $X\ge Y$ and $-1$ otherwise, satisfies
\[
 G(x,y)=\prod_{i=0}^{k-1}E_i
   +\sum_{i=0}^{k-1}\frac{x_i-y_i}{2}
       \prod_{j=i+1}^{k-1}E_j.
\]
The first term handles equality; otherwise the highest differing bit
contributes the comparison sign and all other summands vanish.
Each $E_i$ and $(x_i-y_i)/2$ has Fourier $\ell^1$ norm $1$.
Submultiplicativity and the triangle inequality therefore give
$\mathcal A(G)\le k+1$.  The selector target is
\[
 f_k(x,y,z)=\frac{1-z}{2}G(x,y)-\frac{1+z}{2},
\]
so $\mathcal A(f_k)\le k+2$.  Apply
Theorem~\ref{thm:fourier-teacher} with $d=2k+1$.
\end{proof}

\section{From energy separation to the actual posterior}
\label{post:section}

We now control the posterior at a prescribed width-dependent variance,
rather than first sending the variance to zero separately at each width.
The correction $A_N$ may tend to infinity arbitrarily slowly.
Throughout this section $d\ge1$ is fixed, $D=d+1$, and
\[
 \sigma_N^2=\tau^2N^{-\gamma},\qquad
 \gamma\ge1,\quad \tau>0,\quad A_N\longrightarrow\infty.
\]
All constants in the estimates may depend on the fixed dimension, the
fixed finite training histories, $\gamma$, and $\tau$, but not on $N$.

\begin{proposition}[A direct posterior comparison]
\label{thm:direct}
Let $\mathcal S$ be a fixed finite labeled history. Suppose its affine
unit-margin problem has minimum-norm solution $a_0$, with
$Q_0=\|a_0\|^2>0$, and
\[
 |a_0\cdot\bar x|\ge1\qquad\text{for every }x\in\{-1,1\}^d.
\]
Let $(x_*,y_*)$ be a further labeled point such that
$y_*a_0\cdot\bar x_*<0$. Suppose the affine unit-margin problem for
$\mathcal S\cup\{(x_*,y_*)\}$ is feasible, with squared minimum norm
$Q_1>Q_0$.
Under the prior~\eqref{eq:prior}, conditioned on the global absolute
output margin greater than $1/2$ and on the exact labels in $\mathcal S$,
the posterior probability of label $y_*$ at $x_*$ tends to zero as
$N\to\infty$.
The conclusion holds simultaneously for any fixed finite collection of
such histories and next points.
\end{proposition}
\begin{proof}
Write $A=A_N$ and $\sigma=\sigma_N$. At every sufficiently large width,
Theorem~\ref{thm:scalar} applies to both affine problems, because their
required penalties are fixed and $A_N\to\infty$. Put
\[
 E_0=\Phi_A(Q_0),\qquad E_1=\Phi_A(Q_1),\qquad
 \Delta=Q_1-Q_0>0.
\]
The old problem has a cloned scalar witness of energy $E_0$, with the
same full-cube sign function as $a_0$. Every network fitting the old
history and the additional correct label, with output margins at least
$1/2$ there, has energy at least $E_1$. The latter assertion includes
all nonlinear output functions.

Let $r_A(Q)$ minimize the scalar profile~\eqref{eq:profile}. Its
stationarity equation is
\[
 A=4dQ\,r_A(Q)\frac{\tanh^3r_A(Q)}{\operatorname{sech}^2r_A(Q)}.
\]
The right side increases strictly with its radius. At
$r=(\log A)/4$ it is at most $dQr e^{2r}$, which is smaller than $A$
for all sufficiently large $A$. Thus $r_A(Q_1)\ge(\log A)/4$ eventually.
Evaluating the profile for $Q_0$ at the minimizer for $Q_1$ gives
\begin{equation}\label{post:gap}
 E_1-E_0\ge d\Delta\,r_A(Q_1)^2
 \ge c(\log A)^2,
 \qquad c=d\Delta/16>0.
\end{equation}
Also, for $i=0,1$,
\begin{equation}\label{post:baseline}
 E_i=A/4+O((\log A)^2),
\end{equation}
because the readout contribution is at least $A/4$, and the cloned
choice $r=(\log A)/2$ gives the matching upper estimate.

\paragraph{A strict-margin ball.}
Use the normalized readout variables $b_j=Nv_j$ and whiten the energy by
\[
 U_j=\sqrt{d/N}\,u_j,\qquad
 B_j=\sqrt{A/N}\,b_j.
\]
There are $P=N(d+2)$ coordinates, each independently distributed as
$N(0,\sigma^2)$, and their squared Euclidean norm is $E_A$.
Set
\[
 G=\frac c2(\log A)^2,\qquad
 \varepsilon_N=\frac{G}{16E_0}.
\]
Amplify the scalar witness's readout by $1+\varepsilon_N$.
For large $N$, $\varepsilon_N<1$, its center energy is at most
$E_0+3G/16<E_0+G/4$, and its signed output on every cube input is at
least $(1+\varepsilon_N)/2$. Its global margin slack is therefore
$\varepsilon_N/2$, where
\[
 \varepsilon_N=\Theta((\log A)^2/A).
\]
In whitened coordinates the output is
\[
 F(x)=\frac1{\sqrt{AN}}\sum_{j=1}^N
 B_j\tanh\!\left(\sqrt{N/d}\,U_j\cdot\bar x\right).
\]
Since $|\tanh|\le1$ and $|\tanh'|\le1$,
\[
 \|\nabla F(x)\|
 \le A^{-1/2}+\sqrt{D/(Ad)}\,\|B\|.
\]
On the radius-one ball about the amplified center, this is bounded by a
fixed constant, uniformly over all cube inputs: the center energy is
$O(A)$. Consequently the Euclidean ball of radius
\[
 r_N=A^{-2}
\]
lies inside the same strict global-margin sign cell for every sufficiently
large $N$. Indeed its output perturbations are $O(A^{-2})$, whereas the
available slack is of order $(\log A)^2/A$. Its maximum energy is at most
$E_0+G/2$, because the additional energy is bounded by
$2\sqrt{E_{\rm center}}r_N+r_N^2=O(A^{-3/2})=o(G)$.

The cube of side $2r_N/\sqrt P$ centered at this witness is contained in
the ball. Integrating the minimum Gaussian density on the cube gives
\begin{align}
 \log\Pp(\text{witness cell})
 &\ge-\frac{E_0+G/2}{2\sigma^2}
       +P\log2-2P\log A-\frac P2\log(2\pi P\sigma^2)\nonumber\\
 &\ge-\frac{E_0+G/2}{2\sigma^2}-C_1N(1+\log A).
 \label{post:lower-mass}
\end{align}
The last bound uses
$P\sigma^2=(d+2)\tau^2N^{1-\gamma}\le(d+2)\tau^2$.
Keeping this factor together is essential when $A_N$ diverges slowly.

\paragraph{All correct-label competitors.}
Their prior event is contained in $\{E_A\ge E_1\}$. Since
$E_A/\sigma^2$ has the chi-square distribution with $P$ degrees of freedom,
the exponential-moment bound gives, whenever $E_1>P\sigma^2$,
\begin{align}
 \log\Pp(\text{correct-label event})
 &\le-\frac{E_1}{2\sigma^2}
       +\frac P2\left[1+\log\frac{E_1}{P\sigma^2}\right]\nonumber\\
 &\le-\frac{E_1}{2\sigma^2}
       +C_2N[1+\log A+(\gamma-1)\log N].
 \label{post:upper-mass}
\end{align}
The condition holds eventually by~\eqref{post:baseline} and $A_N\to\infty$.
The posterior denominator contains the witness cell. Since
$E_1-E_0\ge2G$, dividing~\eqref{post:upper-mass} by
\eqref{post:lower-mass} yields
\begin{align}
 \log\Pp(y_*\text{ at }x_*\mid\mathcal S,\text{ global margin})
 &\le-\frac{3G}{4\sigma^2}
       +C_3N[1+\log A+(\gamma-1)\log N]\nonumber\\
 &=-c' N^\gamma(\log A_N)^2
       +C_3N[1+\log A_N+(\gamma-1)\log N]
 \label{post:ratio}
\end{align}
for a fixed $c'>0$.
At $\gamma=1$, the negative term dominates because $\log A_N\to\infty$.
At $\gamma>1$, it also dominates $N\log N$, since
$N^{\gamma-1}(\log A_N)^2/\log N\to\infty$.
This proves the asserted convergence. A finite collection of histories
allows a common width threshold.
\end{proof}

\section{Polynomial IID learning under the same prior}
\label{post:iid-section}

The preceding lower bound concerns an adversarial input order. The same
prior has a polynomial realizable IID guarantee for every affine threshold
target, including the target used in the online construction.

\begin{lemma}[A gap for nonthreshold output functions]
\label{post:nonthreshold-energy}
Let $M=2^d$. If a Boolean function $g$ is not an affine threshold function,
then every width-$N$ tanh network with
$g(x)F(x)\ge1/2$ at all cube inputs has
\[
 E_A\ge\frac{A M^2}{4(M-1)^2}.
\]
\end{lemma}
\begin{proof}
Let $\E_j$ denote uniform averaging over neurons and put
$C=\E_j|b_j|$. For every input $x$, the margin constraint gives
\[
 \E_j\left[|b_j|\{1-\operatorname{sign}(b_j)g(x)
                  \tanh(u_j\cdot\bar x)\}\right]\le C-1/2.
\]
Assign either sign when $b_j=0$. For each neuron there is some input at
which $\operatorname{sign}(b_j)g(x)u_j\cdot\bar x\le0$; otherwise that
oriented affine score would represent $g$ strictly on the whole cube.
Thus at least one of its bracketed deficits is at least one, and every
deficit is nonnegative. Summing the displayed inequalities over all
$M$ inputs gives $C\le M(C-1/2)$, hence
$C\ge M/[2(M-1)]$. Finally
$E_A\ge A\E_j b_j^2\ge AC^2$.
\end{proof}

\begin{theorem}[IID sample complexity]
\label{thm:iid}
Fix an affine threshold target $f:\{-1,1\}^d\to\{-1,1\}$,
any input distribution $\mathcal D$, and $\epsilon,\beta\in(0,1)$.
If
\[
 m\ge\frac4\epsilon\bigl[(d+1)^2\log2+\log(1/\beta)\bigr],
\]
then, with probability at least $1-\beta$ over $m$ IID examples labeled
by $f$, the posterior under~\eqref{eq:prior} assigns probability greater
than $2/3$ to the correct label on a set of $\mathcal D$-probability at
least $1-\epsilon$, at every sufficiently large width.
In particular posterior-label majority has population error at most
$\epsilon$. The width threshold may depend on $f,d,\epsilon,\gamma,\tau$
and the sequence $A_N$, but can be chosen independently of the training
sample and $\mathcal D$.
\end{theorem}
\begin{proof}
Choose a full-cube affine unit-margin implementation of $f$ with squared
norm $Q_f<\infty$. Cloning it with scalar radius $(\log A)/2$ gives
energy $A/4+O((\log A)^2)$. The same small readout amplification and
radius-$A^{-2}$ whitened ball used above give a strict full-target cell
whose prior mass is at least
\[
 \exp\left[-\frac{A/4+O((\log A)^2)}{2\sigma_N^2}
            -C N(1+\log A)\right].
\]
This cell is consistent with every sample labeled by $f$.
By Lemma~\ref{post:nonthreshold-energy}, the union of all nonthreshold
output cells has energy at least
\[
 \frac{AM^2}{4(M-1)^2}\ge\frac A4+\frac A{2M}.
\]
Since $(\log A)^2=o(A)$, applying the same chi-square bound and comparing
with the full-target cell shows that the posterior mass of all
nonthreshold functions is at most a quantity $B_N\to0$, uniformly over
every target-consistent history. Explicitly, for large $N$ one may take
\[
 B_N=\min\left\{1,
 \exp\left[-\frac{A_N}{8M\sigma_N^2}
       +C' N\{1+\log A_N+(\gamma-1)\log N\}\right]\right\}.
\]
The exponent tends to minus infinity for every diverging $A_N$ and
$\gamma\ge1$.

Let $\mathcal H_d$ be the affine threshold functions on the cube.
The $2^d$ input hyperplanes in coefficient space give the standard region
bound
\[
 |\mathcal H_d|\le\sum_{j=0}^{d+1}\binom{2^d}{j}
 \le(2^d+1)^{d+1}\le2^{(d+1)^2}.
\]
A realizable consistency union bound therefore shows that, except with
probability at most $|\mathcal H_d|e^{-m\epsilon/4}\le\beta$, every
threshold function consistent with the sample has population error at
most $\epsilon/4$.
On this sample event, the posterior mean population error is at most
$\epsilon/4+B_N$. Once $B_N\le\epsilon/12$, it is at most $\epsilon/3$.
At an input where the posterior probability of the correct label is at
most $2/3$, the posterior wrong-label probability is at least $1/3$.
Markov's inequality thus bounds the probability of these inputs by
$\epsilon$. This also bounds majority error, and proves the theorem.
\end{proof}


\section{Proof of the main theorem and scope}

The fixed calibration and input order in Lemma~\ref{lem:selector} give,
just before each designated query, a convex affine margin minimum $Q_t$
whose optimizer has absolute margin at least one on the entire cube and
predicts the next label incorrectly. Adding the correct next label raises
the affine minimum to $Q_{t+1}>Q_t$. The exact all-width theorem
\ref{thm:scalar} applies eventually because $A_N\to\infty$.
The posterior comparison in the preceding section therefore sends the
correct-label probability to zero at each step. There are finitely many
steps for each fixed $k$, so one width threshold suffices for all of them.
Using a per-round probability allowance $\eta/T_k$ and a union bound
proves the posterior-sampling assertion without any independence
assumption about predictions. The teacher construction and the IID result
prove the remaining statements.

\begin{proposition}[The scalar witnesses are beaten at the reference scaling]
\label{prop:canonical-comparison}
At $A=1$, for every $N\ge2n_k$, the true target has a strict-margin
implementation of energy at most
\[
 d(d+1)+18(k+2)^2(2d+1)=O(k^3).
\]
Every cloned scalar implementation of the affine witness in
Lemma~\ref{lem:selector} has energy at least
\[
 \sqrt{dQ_t}\ge\sqrt{\frac{2d}{3}(4^k-1)}.
\]
Thus, for all sufficiently large $k$, the true target beats every one of
these scalar witnesses at every such width under the reference scaling.
\end{proposition}
\begin{proof}
The construction in Theorem~\ref{thm:fourier-teacher} supplies a teacher
with $n\le n_k$ neurons, first-row squared norms at most $d+1$, and
readout coefficients of magnitude $B/n$, where
$B\le3(k+2)\sqrt{2d+1}$. Allocate $m_j$ clones to neuron $j$, with
$\sum_jm_j=N$ and each $m_j$ either $\lfloor N/n\rfloor$ or
$\lceil N/n\rceil$. Divide its readout coefficient by $m_j$.
The function is unchanged. Since $N\ge2n$, each $m_j\ge N/(2n)$, so its
canonical energy is at most
\[
 \frac dN\sum_jm_j(d+1)
       +N\sum_j\frac{(B/n)^2}{m_j}
 \le d(d+1)+2B^2.
\]
For a scalar witness, $\tanh r\le r$ gives
\[
 \Phi_1(Q_t)\ge\min_{r>0}\left(dQ_tr^2+\frac1{4r^2}\right)
 =\sqrt{dQ_t}.
\]
The claimed bound on $Q_t$ follows from \eqref{eq:Qt}. The exponential
lower bound eventually exceeds the polynomial upper bound.
\end{proof}

This comparison does not determine which wide wrong functions may still
have low energy. It shows why copying the scalar affine witnesses is not
itself a proof for $A=1$.

The learner uses fixed ordinary tanh, the full Boolean cube, a constant
input coordinate, and a Gaussian prior with exact first-coordinate
isotropy. The argument includes nonlinear competitors at arbitrary width.
The readout correction $A_N\to\infty$ remains essential to this proof:
it makes the minimum-energy hidden units saturate on the active training
constraints, so affine margin geometry determines the optimum. The argument gives no conclusion
for the reference case $A_N\equiv1$, nor for a fixed bounded readout
correction uniformly across this growing target family. Those are separate
questions. The theorem is a statistical statement about an exact posterior;
it supplies no efficient posterior-sampling algorithm.

\end{document}
